
%-------------------------
% Resume in Latex
% Original Author : Sourabh Bajaj
% Adaptation : Hyunggi Chang
% License : MIT
%------------------------

\documentclass[letterpaper,11pt]{article}

\usepackage{latexsym}
\usepackage[empty]{fullpage}
\usepackage{titlesec}
\usepackage{marvosym}
\usepackage[usenames,dvipsnames]{color}
\usepackage{verbatim}
\usepackage{enumitem}
\usepackage[hidelinks]{hyperref}
\usepackage{fancyhdr}
\usepackage[english]{babel}
\usepackage{tabularx}
\usepackage{amsmath}
\usepackage{kotex} % Enable Korean!

\pagestyle{fancy}
\fancyhf{} % clear all header and footer fields
\fancyfoot{}
\renewcommand{\headrulewidth}{0pt}
\renewcommand{\footrulewidth}{0pt}

% Adjust margins
\addtolength{\oddsidemargin}{-0.5in}
\addtolength{\evensidemargin}{-0.5in}
\addtolength{\textwidth}{1in}
\addtolength{\topmargin}{-0.5in}
\addtolength{\textheight}{1.0in}

\urlstyle{same}

\raggedbottom
\raggedright
\setlength{\tabcolsep}{0in}

% Sections formatting
\titleformat{\section}{
  \vspace{-4pt}\scshape\raggedright\large
}{}{0em}{}[\color{black}\titlerule \vspace{-2pt}]

%-------------------------
% Custom commands
\newcommand{\resumeTilde}[0]{
  {\raise.17ex\hbox{$\scriptstyle\mathtt{\sim}$}}
}

\newcommand{\resumeItem}[1]{
  \item\small{
    {#1 \vspace{-2pt}}
  }
}

\newcommand{\resumeSummary}[1]{
  \small{
    \begin{tabular*}{0.97\textwidth}[t]{l@{\extracolsep{\fill}}r}
      #1
    \end{tabular*}
  }
}

\newcommand{\resumeSubheading}[4]{
  \vspace{-1pt}\item
    \begin{tabular*}{0.97\textwidth}[t]{l@{\extracolsep{\fill}}r}
      \textbf{#1} & #2 \\
      \textit{\small#3} & \textit{\small #4} \\
    \end{tabular*}\vspace{-5pt}
}

\newcommand{\resumeEmployment}[4]{
  \vspace{-1pt}\item
    \begin{tabular*}{0.97\textwidth}[t]{l@{\extracolsep{\fill}}r}
      \textbf{#1} & #2 \\
      \textit{\small#3} & \textit{\small #4} \\
    \end{tabular*}\vspace{-5pt}
}

\newcommand{\resumeProject}[2]{
  \vspace{-1pt}\item
    \begin{tabular*}{0.97\textwidth}[t]{l@{\extracolsep{\fill}}r}
      \textbf{#1} \\
      \small{#2} \\
    \end{tabular*}\vspace{-5pt}
}

\newcommand{\resumeResearch}[5]{
  \vspace{-1pt}\item
    \begin{tabular*}{0.97\textwidth}[t]{l@{\extracolsep{\fill}}r}
      \textbf{#1} & #2 \\
      \textit{\small#3} {\small #4 \vspace{-2pt}} & \textit{\small #5} \\
    \end{tabular*}\vspace{-5pt}
}

\newcommand{\resumeTalk}[2]{
  \vspace{-1pt}\item
    \begin{tabular*}{0.97\textwidth}[t]{l@{\extracolsep{\fill}}r}
      \textbf{#1} & #2 \\
    \end{tabular*}\vspace{-5pt}
}

\newcommand{\resumeSkills}[1]{
  \item
    \begin{tabular*}{0.97\textwidth}[t]{l@{\extracolsep{\fill}}r}
      #1
    \end{tabular*}
}

\newcommand{\resumeCommunity}[3]{
  \vspace{-1pt}\item
    \begin{tabular*}{0.97\textwidth}[t]{l@{\extracolsep{\fill}}r}
      \textbf{#1} & #2 \\
      \textit{\small#3} \\
    \end{tabular*}\vspace{-5pt}
}

\newcommand{\resumeSubItem}[2]{\resumeItem{#1}{#2}\vspace{-4pt}}

\renewcommand{\labelitemii}{$\circ$}

\newcommand{\resumeSubHeadingListStart}{\begin{itemize}[leftmargin=*]}
\newcommand{\resumeSubHeadingListEnd}{\end{itemize}}

\newcommand{\resumeEmploymentListStart}{\begin{itemize}[leftmargin=*]}
\newcommand{\resumeEmploymentListEnd}{\end{itemize}}

\newcommand{\resumeItemListStart}{\begin{itemize}}
\newcommand{\resumeItemListEnd}{\end{itemize}\vspace{-5pt}}

%-------------------------------------------
%%%%%%  CV STARTS HERE  %%%%%%%%%%%%%%%%%%%%%%%%%%%%

\begin{document}

%----------HEADING-----------------
\begin{tabular*}{\textwidth}{l@{\extracolsep{\fill}}r}
  \textbf{\href{http://skyoo2003.github.io}{\Large 유성규}} & Github: \href{https://github.com/skyoo2003}{@skyoo2003} $|$ LinkedIn: \href{https://www.linkedin.com/in/skyoo2003}{Sungkyu Yoo} \\
  \href{http://skyoo2003.github.io}{http://skyoo2003.github.io} & Email : \href{mailto:skyoo2003@gmail.com}{skyoo2003@gmail.com} \\
\end{tabular*}

%-----------Summary-----------------
\section{소개}
  \resumeSummary{
    10년 이상 검색 서비스 및 플랫폼을 개발해 온 백엔드 엔지니어입니다. \\
    특히 Elasticsearch와 Kubernetes 기반으로, 시스템 설계부터 운영까지 모두 경험했습니다. \\
    최근에는 LLM/GPU 추론 시스템을 구축하며 비용 효율적인 AI 서비스 개발 역량을 갖추었습니다. \\
    견고하고 확장성 있는 시스템 설계를 통해 비즈니스 성장에 기여하고 싶습니다.
  }

%-----------EXPERIENCE-----------------
\section{경력}
  \resumeEmploymentListStart
    \resumeEmployment
      {\href{https://www.sktelecom.com}{SKTelecom}}{서울 을지로, SKT타워}
      {백엔드 엔지니어 (정규직)}{2022.09 - 2025.10}
        \resumeItemListStart
            \resumeItem{에이닷 및 ICT Family(B tv, Flo 등) 검색서비스 개발/운영}
            \resumeItem{LLM 대화 서비스(B tv 미디어 에이전트) 개발/운영}
            \resumeItem{검색 시스템을 위한 쿠버네티스 클러스터 개발/운영}
            \resumeItem{GPU 인퍼런스 서빙 개발/운영}
        \resumeItemListEnd   
        
    \resumeEmployment
      {\href{https://www.kakaoenterprise.com/}{Kakao Enterprise Corp.}}{경기도 판교, H스퀘어}
      {백엔드 엔지니어 (정규직)}{2019.12 - 2022.08}
      \resumeItemListStart
          \resumeItem{검색서비스 운영 자동화 플랫폼 개발/운영}
          \resumeItem{검색시스템 표준 쿠버네티스 클러스터 개발/운영}
      \resumeItemListEnd

    \resumeEmployment
      {\href{https://www.kakaocorp.com/}{Kakao Corp.}}{경기도 판교, H스퀘어}
      {백엔드 엔지니어 (정규직)}{2015.03 - 2019.11}
        \resumeItemListStart
            \resumeItem{다음 및 카카오톡 검색서비스 개발/운영}
            \resumeItem{검색서비스 운영 자동화 플랫폼 개발/운영}
        \resumeItemListEnd   
  \resumeEmploymentListEnd

%-----------Projects-----------------
\section{프로젝트}
  \resumeSubHeadingListStart
    \resumeProject
      {\textbf{vLLM 기반의 LLM GPU 추론 시스템 구축}}
      {추론 서버와 언어 모델을 최적화하고, 서비스 적용으로 비용을 절감하고 자체 운영 환경 구축}
      \resumeItemListStart
        \resumeItem{추론 서버 및 언어 모델 최적화(양자화)를 검토하고 서비스에 적용}
        \resumeItem{쿠버네티스에 배치하기 위한 리소스를 작성해서 배포하고 프로덕션 적용}
        \resumeItem{\textbf{주요 기여 및 결과}: 뉴스 검색의 요약 Task에서 GPT-4o-mini를 Gemma-3-12B 모델로 대체하여 토큰 비용 절감}
      \resumeItemListEnd

    \resumeProject
      {NVIDIA Triton 기반 임베딩 추론 시스템 구축}
      {사용자 발화를 실시간(Online) 추론을 가능하게 하고, 이를 활용해 벡터 서치 기반으로 검색서비스 품질 고도화}
      \resumeItemListStart
        \resumeItem{서비스에 적용 가능한 Nvidia GPU 추론 서버와 모델 변환(ONNX)을 검토하고 동작 검증}
        \resumeItem{오픈소스 추론 서버를 쿠버네티스에 배치하기 위한 리소스를 작성해서 배포하고 프로덕션 적용}
        \resumeItem{\textbf{주요 기여 및 결과}: 뉴스 검색 시맨틱 필터링으로 저품질 CS 77\% 감소, 발화-청크 유사도 검색 및 키워드 확장에 활용}
      \resumeItemListEnd
        
    \resumeProject
      {On-premise 검색 시스템의 Kubernetes 전환을 통한 구조 개선 및 효율화}
      {검색 서비스를 클라우드 기반으로 구축하여 서비스 개발을 기민하게 하고, 반복적인 작업을 자동화해서 효율적인 운영 달성}
      \resumeItemListStart
        \resumeItem{On-prem 환경에 쿠버네티스 클러스터 구축 및 어플리케이션 배포 자동화(CD) 개발/운영}
        \resumeItem{검색서비스 어플리케이션(ES, FastAPI) 컨테이너화 및 쿠버네티스 배포 개발/운영}
        \resumeItem{쿠버네티스의 컨테이너 메트릭/로그 모니터링 개발/운영}
        \resumeItem{\textbf{주요 기여 및 결과}: 총 4개의 검색서비스를 전환하여 On-prem 대비 53\% 장비 수량 감소로 서버 비용 절감}
      \resumeItemListEnd

    \resumeProject
      {A. 및 ICT Family 검색서비스 개발/운영}
      {서비스 지향 검색 서비스를 개발하고 안정적으로 운영}
      \resumeItemListStart
        \resumeItem{Elasticsearch 구축/운영, 색인 자동화 및 검색 엔진 프론트(검색/랭킹 비즈니스 로직) 개발}
        \resumeItem{검색 서비스 무중단 운영, 모니터링, 지표 개발 및 트러블슈팅}
        \resumeItem{담당 서비스: 에이닷 큐피드(QnA), A./NUGU 유사음악 추천, B tv 영상 추천, B tv 미디어 에이전트 대화/검색}
        \resumeItem{
          \textbf{주요 기여 및 결과}: \\
          * B tv 미디어 에이전트 대화/검색: LLM 응용 서비스 및 ES8 활용 검색 신규 과제 참여. NUGU 음성 검색이 미디어 에이전트로 전환 후 5.57배 트래픽 증가, 만족도 지표 상승.
          * A./NUGU 유사음악 추천: ES8 활용 검색 로직 개편 참여. MAP@K 6.74\%, 트랙 평균 청취 시간 12.74\% 상승 \\
          * B tv 영상 추천: ES7 활용 검색/랭킹 로직 개편 참여. 큐레이션 5.8\%, 유사 추천 3.2\%, 미수급 추천 4\% CTR 상승 \\
        }
      \resumeItemListEnd

    \resumeProject
      {검색시스템 표준 쿠버네티스 프로비저너 설계 및 개발}
      {표준화된 쿠버네티스 클러스터 구축으로 검색 파트 내 쿠버네티스 활용을 확산하고 효율적인 구축 달성}
      \resumeItemListStart
        \resumeItem{쿠버네티스 클러스터 스펙(K8S, Container Runtime/Network 등)을 조사하고, 검색에 적합한 환경 정의하고 자동화}
        \resumeItem{쿠버네티스 어플리케이션 배포에 GitOps 방식을 도입하여 소스코드 버전 관리로 클러스터 형상을 관리}
      \resumeItemListEnd

    \resumeProject
      {검색 서비스 배포/운영 자동화 플랫폼 개발}
      {반복적이고 소모적인 검색서비스 구축 과정을 자동화하여 검색서비스 개발을 효율적으로 할 수 있는 시스템 구축 달성}
      \resumeItemListStart
        \resumeItem{검색 서비스 라이프사이클 자동화 개발}
        \resumeItem{검색 서비스 개발/운영을 위한 백오피스 어드민 개발}
      \resumeItemListEnd

    \resumeProject
      {Daum/Kakao 대용량 검색 서비스 개발 및 운영}
      {다음 포털 통합 검색의 대용량(수백 -- 수천 Tps) 트래픽을 커버할 수 있는 검색서비스 개발과 검색 시스템의 안정적인 운영 달성}
      \resumeItemListStart
        \resumeItem{담당 컬렉션: 다음 게시판 검색, 모바일앱 검색, 책 검색, 요리 검색, 카카오톡 오픈채팅방 검색}
        \resumeItem{검색엔진 색인 자동화, 카카오 자체 검색 엔진 구축, 검색엔진 프론트(검색/랭킹 비즈니스 로직) 서버 개발}
        \resumeItem{검색 서비스 무중단 운영, 모니터링 및 트러블슈팅}
      \resumeItemListEnd
  \resumeSubHeadingListEnd

%-----------Skills-----------------
\section{사용 기술}
  \resumeSubHeadingListStart
    \resumeSkills{\textbf{Languages} - Python, Golang, Java}
    \resumeSkills{\textbf{Frameworks} - FastAPI}
    \resumeSkills{\textbf{DevOps} - Kubernetes, Helm, Kustomize, Ansible}
    \resumeSkills{\textbf{Solutions} - Elasticsearch, Prometheus/Thanos, ArgoCD}
    \resumeSkills{\textbf{OS} - RHEL, Ubuntu}
  \resumeSubHeadingListEnd

%-----------EDUCATION-----------------
\section{교육}
  \resumeSubHeadingListStart
    \resumeSubheading
      {한국외국어대학교} {경기도 용인}
      {컴퓨터공학과 학사 (4.23/4.5)} {2009.03 -- 2015.02}
      \resumeItemListStart
        \resumeItem{지능형 문서 보유목관리 시스템 개발 - (주)STORYANT, Part-time}
        \resumeItem{SDR 운용 소프트웨어 품질 평가 지표 및 평가 시스템 개발 - (주)랏데커뮤니케이션, Part-time}
        \resumeItem{학교자치기구(Hufstory) 활동 - 강의평가 웹 앱(Hubigo), 교내 셔틀 버스 위치 추적 앱 개발}
        \resumeItem{대학생 알고리즘 프로그래밍 경진대회(ACM-ICPC) 2013 예선, 2014 지역 본선 참가}
      \resumeItemListEnd
  \resumeSubHeadingListEnd

%-----------Talks-----------------
\section{외부 활동}
    \resumeSubHeadingListStart
      \resumeTalk{\href{https://devocean.sk.com/blog/techBoardDetail.do?ID=167143}{쿠버네티스를 활용한 검색추천 시스템 구축과 미디어 에이전트 적용 사례}}{2024.12}
        \resumeItemListStart
          \resumeItem{쿠버네티스 클러스터 구축 과정을 정리하고, 미디어 에이전트 서비스 적용 사례를 기술 블로그로 작성}
        \resumeItemListEnd
      \resumeTalk{\href{https://docs.google.com/document/d/1zqta43b_VOd1f9PsCzUMJIy9P75sSaWqf3Hkr9IlpSc}{etcd deep-dive 스터디}}{2023.03 -- 2023.05}
        \resumeItemListStart
          \resumeItem{쿠버네티스 상태 관리의 핵심인 etcd를 깊게 이해하고 업무에 활용하기 위한 스터디 참여}
        \resumeItemListEnd
    \resumeSubHeadingListEnd

%-------------------------
% Resume in Latex
% Original author : Sourabh Bajaj
% Adaptation : Hyunggi Chang (changh95)
% License : MIT
%------------------------

%-------------------------
% Custom commands - Do Look into this area, if you wish to customise further!
% 포맷이 마음에 들지 않으시다면 이 부분을 수정하세요!

\newcommand{\resumeSubProject}[2]{
  \item{
    {#1 \vspace{0pt}}
    {#2}
  }
}

\newcommand{\resumeProjectListStart}{\begin{itemize}[leftmargin=*]}
\newcommand{\resumeProjectListEnd}{\end{itemize}}

\newcommand{\resumeSubProjectListStart}{\begin{itemize}[leftmargin=*]}
\newcommand{\resumeSubProjectListEnd}{\end{itemize}}

%-------------------------------------------
%%%%%%  CV STARTS HERE  %%%%%%%%%%%%%%%%%%%%%%%%%%%%

\newpage

%----------HEADING-----------------
\begin{tabular*}{\textwidth}{l@{\extracolsep{\fill}}r}
  \textbf{\href{http://skyoo2003.github.io}{\Large 유성규 (경력 기술서)}} & Github: \href{https://github.com/skyoo2003}{@skyoo2003} $|$ LinkedIn: \href{https://www.linkedin.com/in/skyoo2003}{Sungkyu Yoo} \\
  \href{http://skyoo2003.github.io}{http://skyoo2003.github.io} & Email : \href{mailto:skyoo2003@gmail.com}{skyoo2003@gmail.com} \\
\end{tabular*}

%-----------PROJECTS-----------------
\section{\textbf{vLLM 기반의 LLM GPU 추론 시스템 구축}}
  \resumeProjectListStart
    \resumeProject
    {\textbf{소속}: SKTelecom} {2025.06 - 2025.10}

    {
      %----------- [프로젝트 개요] -----------%
      \textbf{프로젝트 개요:}
      \begin{itemize}[leftmargin=*]
          \item 외부 LLM API 사용으로 인한 \textbf{토큰 비용 급증 문제}를 해결하고, 팀 유휴 GPU 자산(A100)의 활용도를 극대화하기 위해 \textbf{자체 LLM 추론 플랫폼 구축 프로젝트를 직접 발제하고 주도}했습니다.
          \item 목표는 LLM 토큰 비용을 절감하기 위해 오픈소스 LLM을 안정적으로 서빙할 수 있는 고성능 추론 환경을 구축하는 것이었습니다.
      \end{itemize}

      %----------- [주요 역할 및 기여] -----------%
      \textbf{주요 역할 및 기여:}
      \begin{itemize}[leftmargin=*]
          \item \textbf{데이터 기반의 최적 추론 서버 스택 선정 및 아키텍처 설계}
              \begin{itemize}[leftmargin=*, label=$\circ$]
                  \item Throughput 및 Latency 지표 분석 결과, 가장 우수한 성능을 보인 \textbf{vLLM을 추론 엔진으로 선정}하고, 이를 Kubernetes 환경에 배포하기 위한 명세 개발을 진행했습니다.
              \end{itemize}

          \item \textbf{모델 양자화(Quantization)를 통한 VRAM 사용량 및 추론 성능 최적화}
              \begin{itemize}[leftmargin=*, label=$\circ$]
                  \item 한정된 GPU(A100) VRAM 내에서 더 많은 동시 요청을 처리하기 위해 AWQ, GPTQ 등 다양한 양자화 기법을 심도 있게 분석하고 테스트했습니다.
                  \item 품질 저하를 최소화하면서도, 실제 서비스 부하 환경에서 \textbf{가장 우수한 Latency를 보인 GPTQ 양자화 기법을 최종 도입}하여 VRAM 사용량을 크게 줄이고 추론 속도를 향상시켰습니다.
              \end{itemize}
              
          \item \textbf{Kubernetes 기반의 안정적인 추론 서비스 배포 및 모니터링}
              \begin{itemize}[leftmargin=*, label=$\circ$]
                  \item 구축된 추론 서버를 컨테이너화하고, GPU 자원을 효율적으로 관리할 수 있도록 Kubernetes 클러스터에 배포하는 CI/CD 파이프라인을 구축했습니다.
                  \item Prometheus/Grafana를 통해 GPU 활용률, 추론 속도, 처리량 등 성능 지표를 실시간으로 모니터링하여 안정적인 운영 기반을 마련했습니다.
              \end{itemize}
      \end{itemize}

      %----------- [결과 및 성과] -----------%
      \textbf{결과 및 성과:}
      \begin{itemize}[leftmargin=*]
          \item \textbf{고비용의 외부 LLM을 대체할 수 있는 플랫폼을 성공적으로 구축}하여, LLM 토큰 비용 부담을 절감할 수 있는 기반을 마련했습니다.
          \item GPTQ 양자화 모델(Gemma-3-12B) 서빙 시, 원본 모델 대비 품질 저하 없이 \textbf{99th percentile 기준, 토큰 당 생성 시간을 약 20\% 단축}하는 성능을 달성했습니다.
          \item 구축된 플랫폼을 통해 뉴스 요약 Task에서 OpenAI gpt-4o-mini를 gemma-3-12b 모델로 전환하고 내부 처리가 가능하게 되어, \textbf'최적화 가능성과 비용 절감'에 기여했습니다.
      \end{itemize}

      %----------- [사용한 기술] -----------%
      \textbf{사용한 기술:} Python, vLLM, LLM Quantization(AWQ, GPTQ), Kubernetes, GPU(A100), Prometheus, Grafana
    }

  \resumeProjectListEnd

\section{\textbf{NVIDIA Triton 기반 임베딩 추론 시스템 구축}}
  \resumeProjectListStart
    \resumeProject
    {\textbf{소속}: SKTelecom} {2024.06 - 2025.10}

    {
      %----------- [프로젝트 개요] -----------%
      \textbf{프로젝트 개요:}
      \begin{itemize}[leftmargin=*]
          \item 이를 위해, 사용자 질의의 의미를 이해하는 시맨틱 검색(Semantic Search)을 도입하고자 했으며, \textbf{실시간(Online) 임베딩 모델 추론을 위한 고성능 GPU 서빙 시스템을 구축}했습니다.
      \end{itemize}

      %----------- [주요 역할 및 기여] -----------%
      \textbf{주요 역할 및 기여:}
      \begin{itemize}[leftmargin=*]
          \item \textbf{NVIDIA Triton Inference Server 도입을 통한 추론 성능 극대화}
              \begin{itemize}[leftmargin=*, label=$\circ$]
                  \item 상용 트래픽 처리 및 GPU 활용을 극대화하기 위해 전문 추론 서버인 \textbf{NVIDIA Triton을 검토 후 도입}했습니다.
                  \item Triton의 기능인 \textbf{동적 배치(Dynamic Batching)를 활용}하여, 실시간으로 들어오는 단건 요청들을 GPU에서 효율적으로 일괄 처리함으로써 \textbf{Throughput을 향상}시켰습니다.
              \end{itemize}

          \item \textbf{ONNX 모델 변환 및 Kubernetes 기반의 안정적인 배포}
              \begin{itemize}[leftmargin=*, label=$\circ$]
                  \item PyTorch 모델을 \textbf{ONNX(Open Neural Network Exchange) 형식으로 변환}하여 추론 속도를 개선하고 특정 프레임워크에 대한 종속성을 제거했습니다.
                  \item \textbf{Kubernetes의 Nvidia Device Plugin을 활용}하여 GPU 자원을 할당하고 안정적인 서빙 환경을 구축했습니다.
              \end{itemize}
              
          \item \textbf{Prometheus/Grafana 기반의 통합 추론 모니터링 시스템 구축}
              \begin{itemize}[leftmargin=*, label=$\circ$]
                  \item 시스템의 안정적인 운영을 위해, \textbf{GPU 활용률, 메모리 등 인프라 지표}부터 \textbf{추론 지연 시간(Latency), 처리량(Throughput), 에러율} 등 애플리케이션 지표까지 포괄하는 통합 모니터링 대시보드를 구축했습니다.
                  \item 이를 통해 잠재적인 \textbf{성능 병목 구간을 사전에 식별}하고, 장애 발생 시 신속하게 원인을 분석할 수 있는 모니터링 환경을 마련했습니다.
              \end{itemize}
              
          \item \textbf{철저한 성능 검증을 통한 서비스 출시 안정성 확보}
              \begin{itemize}[leftmargin=*, label=$\circ$]
                  \item 서비스 오픈에 앞서, 실제 사용자 로그를 기반으로 한 \textbf{시나리오 기반 부하 테스트를 직접 설계하고 수행}했습니다.
                  \item 목표 TPS 하에서 \textbf{99th percentile 응답 속도가 서비스 SLO(Service Level Objective)를 충족함을 검증}하여, 프로덕션 환경에서의 성능과 안정성을 사전에 확보했습니다.
              \end{itemize}
      \end{itemize}

      %----------- [결과 및 성과] -----------%
      \textbf{결과 및 성과:}
      \begin{itemize}[leftmargin=*]
          \item \textbf{뉴스 검색 관련 저품질 고객 불만(CS) 77\% 감소}라는 검색 품질 개선 성과를 달성했습니다.
          \item \textbf{시맨틱 검색을 통한 검색 품질 개선 기반을 마련}했으며, 구축된 추론 시스템은 '발화-청크 유사도 검색', '키워드 확장' 등 여러 서비스에 다양한 목적으로 재사용되며 기술적 자산이 되었습니다.
      \end{itemize}

      %----------- [사용한 기술] -----------%
      \textbf{사용한 기술:} Python, Kubernetes, NVIDIA Triton Inference Server, ONNX, GPU(A40), Prometheus, Grafana
    }

    \resumeProjectListEnd

  \section{\textbf{On-premise 검색 시스템의 Kubernetes 전환을 통한 구조 개선 및 효율화}}
    \resumeProjectListStart
      \resumeProject
      {\textbf{소속}: SKTelecom} {2023.03 - 2025.10}
      
      {
        %----------- [프로젝트 개요] -----------%
        \textbf{프로젝트 개요:}
        \begin{itemize}[leftmargin=*]
            \item 물리 장비(PM) 기반의 수동 배포 방식으로 운영되던 레거시 검색 인프라의 \textbf{배포 속도, 비효율적인 자원 사용, 운영 복잡성} 문제를 해결하기 위해, Kubernetes 기반의 컨테이너 플랫폼으로 전환하는 프로젝트를 \textbf{End-to-End로 직접 설계하고 구축}했습니다.
            \item 이를 통해 개발 생산성을 극대화하고, 인프라 운영 비용을 획기적으로 절감하는 것을 목표로 했습니다.
        \end{itemize}

        %----------- [주요 역할 및 기여] -----------%
        \textbf{주요 역할 및 기여:}
        \begin{itemize}[leftmargin=*]
            \item \textbf{On-premise 환경에 최적화된 Kubernetes 클러스터 설계 및 구축 주도}
                \begin{itemize}[leftmargin=*, label=$\circ$]
                    \item 자동화 도구(\textbf{Kubespray})를 활용하여 On-premise 환경에 표준화된 Kubernetes 클러스터를 신속하게 배포하는 프로세스를 정립했습니다.
                    \item 네트워크 정책 적용을 위해 \textbf{Calico(CNI)}를, 서비스 노출을 위해 \textbf{Nginx(Ingress Controller)}를 도입하는 등, 컨테이너 기반으로 안정적인 플랫폼을 마련했습니다.
                \end{itemize}

            \item \textbf{ArgoCD를 활용한 주기적인 배포(CD) 구축}
                \begin{itemize}[leftmargin=*, label=$\circ$]
                    \item 클러스터 관리자가 Gitlab에 소스코드 버전 관리로 클러스터 형상을 관리하고 ArgoCD를 활용해서 배포가 가능하도록 구축했습니다.
                    \item 이를 통해 휴먼 에러를 최소화하며, 모든 변경 사항을 코드로 관리할 수 있게 되었습니다.
                \end{itemize}
                
            \item \textbf{검색 서비스의 성공적인 컨테이너화 및 마이그레이션}
                \begin{itemize}[leftmargin=*, label=$\circ$]
                    \item \textbf{Elasticsearch 클러스터}와 \textbf{FastAPI 기반 검색 API 서버}를 컨테이너화하고, Kubernetes 환경으로 성공적으로 이전했습니다.
                    \item Kubernetes를 활용하여 Elasticsearch와 같은 상태 저장 애플리케이션을 안정적으로 운영하는 노하우를 확보했습니다.
                \end{itemize}
                
            \item \textbf{Prometheus/Elastic 기반의 모니터링 시스템 구축}
                \begin{itemize}[leftmargin=*, label=$\circ$]
                    \item 클러스터의 메트릭부터 배포된 애플리케이션의 상태까지 한눈에 파악할 수 있는 통합 대시보드를 구축하여, \textbf{서비스 가시성을 확보}하고 원인 파악이 가능한 환경을 마련했습니다.
                \end{itemize}
        \end{itemize}

        %----------- [결과 및 성과] -----------%
        \textbf{결과 및 성과:}
        \begin{itemize}[leftmargin=*]
            \item \textbf{서버 운영 비용 53\% 절감:} 컨테이너 집적도를 높여 기존 물리 서버 방식 대신 자원 사용률을 획기적으로 개선했습니다.
            \item \textbf{배포 리드타임 단축:} 신규 서비스 배포 및 업데이트에 소요되던 시간을 극적으로 단축(수일 → 수십 분)하여 서비스 개발 기민성을 향상시켰습니다.
            \item \textbf{운영 효율성 증대:} 인프라 구성 및 배포 자동화를 통해 반복적인 수작업을 제거하고, 개발자가 서비스 개발에만 집중할 수 있는 환경을 제공했습니다.
        \end{itemize}

        %----------- [사용한 기술] -----------%
        \textbf{사용한 기술:} Kubernetes, Kubespray, Ansible, ArgoCD, Helm, Elasticsearch, FastAPI, Prometheus, Elastic/Filebeat/Kibana
      }
  \resumeProjectListEnd

\section{\textbf{A. 및 ICT Family 검색서비스 개발/운영}}
  \resumeProjectListStart
    \resumeProject
    {\textbf{소속}: SKTelecom} {2022.09 - 2025.10}

    {
      \resumeSubProjectListStart
        \resumeSubProject{\textbf{LLM과 RAG(검색 증강 생성)를 활용한 대화형 미디어 검색 에이전트 개발}}

        {
          %----------- [프로젝트 개요] -----------%
          \textbf{프로젝트 개요:}
          \begin{itemize}[leftmargin=*]
              \item LLM을 활용하여 기존의 음성 검색을 대체하는 \textbf{대화형 콘텐츠 탐색 서비스}를 구축하는 프로젝트를 담당했습니다.
              \item LLM의 품질과 안정성을 상용 서비스 수준으로 끌어올리기 위해, \textbf{RAG(검색 증강 생성) 패턴을 설계 및 도입}하고 대화 검색 서비스를 개발했습니다.
          \end{itemize}

          %----------- [주요 역할 및 기여] -----------%
          \textbf{주요 역할 및 기여:}
          \begin{itemize}[leftmargin=*]
              \item \textbf{LLM 기반 대화 처리 아키텍처 설계 및 개발}
                  \begin{itemize}[leftmargin=*, label=$\circ$]
                      \item 사용자 발화부터 대화 이력 관리, 의도/개체명 추출(NLU), 검색 연동, 최종 답변 생성에 이르는 \textbf{전체 대화 흐름의 로직을 주도적으로 개발}했습니다.
                  \end{itemize}

              \item \textbf{RAG(검색 증강 생성) 패턴을 선도적으로 도입하여 LLM 한계 극복}
                  \begin{itemize}[leftmargin=*, label=$\circ$]
                      \item 초기 LLM(GPT 3.5)의 문맥 이해도 및 환각(Hallucination) 문제를 해결하기 위해, 사용자 발화와 연관된 VOD 메타데이터를 Elasticsearch에서 실시간으로 검색하여 \textbf{LLM 프롬프트에 동적으로 주입하는 RAG 로직을 직접 구현}했습니다.
                      \item 이를 통해 Few-shot prompting의 한계를 극복하고, B tv 콘텐츠에 대한 정확하고 풍부한 답변 생성을 가능하게 했습니다.
                  \end{itemize}
                  
              \item \textbf{안정적인 서비스 운영을 위한 프롬프트 엔지니어링 및 운영 시스템 구축}
                  \begin{itemize}[leftmargin=*, label=$\circ$]
                      \item LLM이 비즈니스 요구사항에 따라 일관된 결과(의도, 개체명)를 반환하도록 \textbf{정교한 프롬프트를 설계}에 참여 하고, 구체적인 예시를 제공하여 응답 품질을 제어했습니다.
                      \item LLM의 통계적 한계로 인한 돌발적인 오작동에 대응하기 위해, 특정 발화 패턴에 대해서는 사전 정의된 결과를 반환하는 \textbf{운영 기제를 설계}하여 서비스 안정성을 확보했습니다.
                  \end{itemize}
                  
              \item \textbf{검색-LLM 연동 및 성능 최적화}
                  \begin{itemize}[leftmargin=*, label=$\circ$]
                      \item LLM이 추출한 엔티티를 활용해서 Elasticsearch에서 최적의 검색 결과를 도출하는 로직 개발에 참여했습니다.
                      \item LLM의 답변 생성 시간을 고려, 서비스 SLO를 충족시키기 위해 \textbf{최대 출력 토큰 수를 제어하면서도 문장의 완결성을 유지}하는 답변 생성 로직을 개발했습니다.
                  \end{itemize}
          \end{itemize}

          %----------- [결과 및 성과] -----------%
          \textbf{결과 및 성과:}
          \begin{itemize}[leftmargin=*]
              \item 기존 NUGU 음성 검색을 성공적으로 대체하고, B tv 핵심 기능으로 자리매김하여 \textbf{사용자 트래픽 5.57배 증가}라는 성장에 기여했습니다.
              \item LLM 응용 서비스를 성공적으로 상용화한 경험을 통해, 사용자의 만족도를 크게 향상시키고 비즈니스 성장에 직접적으로 기여했습니다.
          \end{itemize}

          %----------- [사용한 기술] -----------%
          \textbf{사용한 기술:} Python(FastAPI), LLM(Prompt Engineering), Elasticsearch 8(RAG), Kubernetes
        }
      \resumeSubProjectListEnd
    
      \resumeSubProjectListStart
        \resumeSubProject{\textbf{ES8 벡터 검색 도입을 통한 유사음악 추천 개선}}

        {
         %----------- [프로젝트 개요] -----------%
          \textbf{프로젝트 개요:}
          \begin{itemize}[leftmargin=*]
              \item 유지보수만 가능했던 레거시 유사음악 추천 시스템을 개선하는 프로젝트를 주도했습니다.
              \item 특히, 별도의 벡터 DB 없이 \textbf{Elasticsearch 8의 Approximate kNN(ANN) 기능을 도입}하여, 시스템 복잡도는 낮추고 검색 품질과 사용자 경험 지표는 극대화하는 것을 목표로 했습니다.
          \end{itemize}

          %----------- [주요 역할 및 기여] -----------%
          \textbf{주요 역할 및 기여:}
          \begin{itemize}[leftmargin=*]
              \item \textbf{Elasticsearch 8 벡터 검색 기술 검증 및 아키텍처 설계 주도}
                  \begin{itemize}[leftmargin=*, label=$\circ$]
                      \item 프로젝트 초기에 Milvus, FAISS 등 외부 벡터 솔루션과 ES8의 내장 ANN 기능 간의 \textbf{기술적 트레이드오프를 분석}했습니다.
                      \item 수십만 건 규모의 문서 수와 운영 효율성을 고려, 별도 시스템 도입의 오버헤드 없이 \textbf{단일 Elasticsearch 클러스터로 데이터 검색과 벡터 검색을 통합하는 아키텍처를 설계하고 도입을 결정}했습니다.
                  \end{itemize}

              \item \textbf{Kubernetes 기반의 안정적인 ES8 클러스터 구축 및 운영}
                  \begin{itemize}[leftmargin=*, label=$\circ$]
                      \item 팀 내 쿠버네티스 환경 위에 벡터 검색 워크로드에 최적화된 리소스(CPU/Memory)를 산정하여 \textbf{ES8 클러스터를 안정적으로 구축}하고, 배포 자동화를 통해 수 분 내에 구성 가능하도록 환경을 마련했습니다.
                  \end{itemize}
                  
              \item \textbf{임베딩 벡터 데이터 처리 파이프라인 개발 및 최적화}
                  \begin{itemize}[leftmargin=*, label=$\circ$]
                      \item HDFS에 저장된 임베딩 벡터와 트랙 메타데이터를 결합하는 \textbf{색인 파이프라인을 개발}했습니다.
                  \end{itemize}

              \item \textbf{추천 API 개발 및 서비스 안정성 확보}
                  \begin{itemize}[leftmargin=*, label=$\circ$]
                      \item FastAPI 기반의 추천 API를 개발하고, 실제 트래픽 기반의 부하 테스트를 통해 \textbf{200 TPS 환경에서 99th percentile 응답속도 204ms를 달성}하여 서비스 SLO를 충족했습니다.
                  \end{itemize}
          \end{itemize}

          %----------- [결과 및 성과] -----------%
          \textbf{결과 및 성과:}
          \begin{itemize}[leftmargin=*]
              \item \textbf{추천 정확도 지표(MAP@10) 6.74\% 상승}
              \item \textbf{사용자 만족도 지표(곡 평균 청취 시간) 12.74\% 상승}
              \item 레거시 시스템을 성공적으로 개편하고, Elasticsearch 내에서 벡터 검색을 구현하는 기술적 선례를 확보하여 팀의 기술 역량을 한 단계 높이는 데 기여했습니다.
          \end{itemize}

          %----------- [사용한 기술] -----------%
          \textbf{사용한 기술:} Python(FastAPI), Elasticsearch 8 (Approximate kNN), Kubernetes, Ansible
        }
      \resumeSubProjectListEnd

      \resumeSubProjectListStart
        \resumeSubProject {\textbf{B tv 추천서비스 개발 및 품질 고도화}}

        {
          %----------- [프로젝트 개요] -----------%
          \textbf{프로젝트 개요:}
          \begin{itemize}[leftmargin=*]
              \item 이원화된 레거시 추천 시스템을 Elasticsearch 기반의 검색 시스템으로 통합하고, 최신 랭킹/추천 모델을 적용하여 \textbf{만족도 지표(CTR)를 개선}하는 것을 목표로 한 프로젝트를 진행했습니다.
              \item 서비스 개발 담당자로서 \textbf{색인 파이프라인부터 API 개발, 성능 최적화, 안정적인 운영}까지 End-to-End 개발을 책임졌습니다.
          \end{itemize}

          %----------- [주요 역할 및 기여] -----------%
          \textbf{주요 역할 및 기여:}
          \begin{itemize}[leftmargin=*]
              \item \textbf{검색엔진 색인 파이프라인 설계 및 구축 주도}
                  \begin{itemize}[leftmargin=*, label=$\circ$]
                      \item VOD 메타 데이터, 지식그래프, 유저 피드백(클릭/노출) 등 다양한 데이터를 결합하여 랭킹 모델에 최적화된 \textbf{Elasticsearch 색인 스키마를 설계}하고, 색인 파이프라인을 개발했습니다.
                      \item 검색 랭킹 및 추천 모델 개발 담당자와 긴밀한 협업을 통해 추천 서빙에 필요한 피처를 정의하고 색인 필드에 반영하여 추천 품질 극대화에 기여했습니다.
                  \end{itemize}

              \item \textbf{Blue/Green 배포 기반의 무중단 색인 자동화 시스템 개발}
                  \begin{itemize}[leftmargin=*, label=$\circ$]
                      \item Python과 Ansible을 활용하여 정적/증분 색인, 리그레션 테스트, Active-Standby 엔진 교체 등 전체 워크플로우를 자동화하여 \textbf{리스크를 최소화}하고 운영 안정성을 확보했습니다.
                      \item 이를 통해 색인 로직 변경에도 서비스 중단 없이 안정적으로 검색 품질을 갱신할 수 있는 기반을 마련했습니다.
                  \end{itemize}
                  
              \item \textbf{벡터 검색을 결합한 추천 API 개발 및 성능 최적화}
                  \begin{itemize}[leftmargin=*, label=$\circ$]
                      \item FastAPI 기반의 추천 API(유사/미수급/큐레이션)를 개발했습니다.
                      \item ES7 기반에서  '유사 콘텐츠 추천' 기능을 구현하기 위해 \textbf{FAISS 벡터 검색 라이브러리를 도입}하여 검색 시스템을 구성했습니다.
                      \item 실사용자 로그 기반으로 부하 테스트를 수행하고 \textbf{200 TPS 환경에서 99백분위 응답시간을 큐레이션 54ms, 유사 추천 296ms를 달성}하여 서비스 SLO를 충족했습니다.
                  \end{itemize}
          \end{itemize}

          %----------- [결과 및 성과] -----------%
          \textbf{결과 및 성과:}
          \begin{itemize}[leftmargin=*]
              \item \textbf{만족도 지표(CTR) 상승 달성}: 큐레이션 5.8\%, 유사 추천 3.2\%, 미수급 추천 4.0\% CTR 지표 상승
              \item 레거시 시스템을 성공적으로 개편하고, 서비스 SLO 내에서 안정적인 서비스 운영이 가능하도록 프로젝트를 개발했습니다.
          \end{itemize}

          %----------- [사용한 기술] -----------%
          \textbf{사용한 기술:} Python(w/ FastAPI), Elasticsearch 7, FAISS, Ansible

        }
      \resumeSubProjectListEnd
    }

  \resumeProjectListEnd

\section{\textbf{검색시스템 표준 쿠버네티스 프로비저너 설계 및 개발}}
  \resumeProjectListStart
    \resumeProject
    {\textbf{소속}: Kakao Enterprise Corp.} {2019.12 - 2022.08}

    {
      %----------- [프로젝트 개요] -----------%
      \textbf{프로젝트 개요:}
      \begin{itemize}[leftmargin=*]
          \item 팀별로 파편화되어 운영되던 쿠버네티스 클러스터로 인해 발생하는 \textbf{운영 리소스 증가, 장애 대응 지연, 기술 노하우 단절} 등의 문제를 해결하기 위해, 표준 쿠버네티스를 구축하는 프로젝트에 참여했습니다.
          \item 목표는 수 분 내에 표준화된 클러스터를 '구축'하고, \textbf{GitOps 기반으로 모든 변경 이력을 투명하게 관리}하여 운영 안정성을 획기적으로 높이는 자동화를 구현하는 것이었습니다.
      \end{itemize}

      %----------- [주요 역할 및 기여] -----------%
      \textbf{주요 역할 및 기여:}
      \begin{itemize}[leftmargin=*]
          \item \textbf{Kubespray 기반의 표준화된 클러스터 구축 자동화}
              \begin{itemize}[leftmargin=*, label=$\circ$]
                  \item 오픈소스 \textbf{Kubespray}의 내부 구조를 분석하여, 검색 시스템에 최적화된 컨테이너 런타임(containerd), 네트워크(flannel) 등을 포함한 \textbf{표준 클러스터 형상을 정의}하고 Ansible 기반의 배포 스크립트를 개발했습니다.
                  \item 인증서 갱신, 노드 레이블링 등 클러스터 운영에 필수적인 반복 작업을 자동화하여 수동 작업으로 인한 휴먼 에러를 원천적으로 차단했습니다.
              \end{itemize}

          \item \textbf{ArgoCD를 활용한 GitOps 기반 클러스터 형상 관리 체계 구축}
              \begin{itemize}[leftmargin=*, label=$\circ$]
                  \item YAML 파일 수동 관리 시 발생하던 \textbf{의도치 않은 설정 롤백 및 배포 충돌 문제}를 해결하기 위해, \textbf{ArgoCD를 도입하여 GitOps 워크플로우를 주도적으로 설계}하고 구축했습니다.
                  \item 'Gitlab Flow' 기반의 브랜치 전략을 도입, 코드 리뷰를 통해 모든 변경 사항을 검증하고, Git을 통해 클러스터의 모든 상태를 추적/관리함으로써 \textbf{배포 안정성과 가시성을 획기적으로 향상}시켰습니다.
              \end{itemize}
              
          \item \textbf{자동화된 클러스터 검증 프로세스 개발}
              \begin{itemize}[leftmargin=*, label=$\circ$]
                  \item 서버 입고 시 스펙 불일치, 설정 누락 등 사소하지만 치명적인 문제를 방지하기 위해, 클러스터 구축 전후의 \textbf{서버 상태를 자동으로 검증하는 체크리스트를 Ansible으로 개발}하여 전체 파이프라인에 통합했습니다.
              \end{itemize}
      \end{itemize}

      %----------- [결과 및 성과] -----------%
      \textbf{결과 및 성과:}
      \begin{itemize}[leftmargin=*]
          \item \textbf{다수의 프로젝트의 쿠버네티스 인프라 기반으로 채택 및 기술 확산}
              \begin{itemize}[leftmargin=*, label=$\circ$]
                  \item \textbf{썸네일 캐시 저장소:} 이미지 검색에서 활용할 썸네일 캐시를 분산 키-값 저장소 구축에 활용.
                  \item \textbf{MLOps 플랫폼:} GPU 노드를 포함한 멀티 클러스터 환경을 본 플랫폼으로 구축.
                  \item \textbf{API Gateway:} 수십 대 규모의 게이트웨이 클러스터를 안정적으로 구축 및 운영.
                  \item \textbf{Search Cloud:} 테넌트별로 분리된 수백 대 규모의 클러스터 구축 기반으로 활용.
              \end{itemize}
          \item \textbf{클러스터 구축 시간을 짧게는 수시간에서 수일 내로 단축}하고, 표준화된 운영 방식을 전파하여 조직의 전체적인 쿠버네티스 기술 역량 향상에 기여했습니다.
      \end{itemize}

      %----------- [사용한 기술] -----------%
      \textbf{사용한 기술:} Kubernetes, Ansible(Kubespray), ArgoCD, GitOps
    }

  \resumeProjectListEnd

\section{\textbf{검색 서비스 배포/운영 자동화 플랫폼 개발}}
  \resumeProjectListStart
    \resumeProject
    {\textbf{소속}: Kakao Corp. / Kakao Enterprise Corp.} {2019.12 - 2022.08}
    
    {
      %----------- [프로젝트 개요] -----------%
      \textbf{프로젝트 개요:}
      \begin{itemize}[leftmargin=*]
          \item 수십 개에 달하는 검색 서비스의 반복적인 배포/운영 업무가 개발팀의 기능 개발을 저해하는 문제를 해결하기 위해, \textbf{중앙화된 자동화 플랫폼을 구축}하는 프로젝트에 참여했습니다.
          \item 목표는 검색 서비스의 전체 라이프사이클(서버 구성, 배포, 색인, 모니터링)을 자동화하고, 검색시스템을 구축하고 운영할 수 있는 환경을 제공하여 \textbf{조직 전체의 개발 생산성을 극대화}하는 것이었습니다.
      \end{itemize}

      %----------- [주요 역할 및 기여] -----------%
      \textbf{주요 역할 및 기여:}
      \begin{itemize}[leftmargin=*]
          \item \textbf{IaC(Infrastructure as Code) 기반의 자동화 워크플로우 설계}
              \begin{itemize}[leftmargin=*, label=$\circ$]
                  \item \textbf{Ansible}을 활용하여 서버 프로비저닝, 검색엔진 클러스터 구성, API 서버 배포 등 인프라 구성을 포함한 운영 작업을 코드로 정의하고 형상 관리했습니다.
                  \item 워크플로우 엔진(\textbf{Digdag})을 도입하여, '정적/증분 색인', 'Blue/Green 배포 및 롤백' 등 운영 시나리오를 실행하는 자동화 파이프라인을 개발에 참여했습니다.
              \end{itemize}
 
          \item \textbf{검색 서비스 운영 표준화 정립}
              \begin{itemize}[leftmargin=*, label=$\circ$]
                  \item 서비스별로 파편화되어 있던 컴포넌트 구성, 디렉토리 구조, 설정 관리 방식을 표준화하여 플랫폼 기반의 자동화가 효율적으로 동작할 수 있는 기반을 마련했습니다.
              \end{itemize}

          \item \textbf{검색서비스 개발자를 위한 백오피스 어드민(플랫폼 포털) 개발}
              \begin{itemize}[leftmargin=*, label=$\circ$]
                  \item 개발자들이 UI를 통해 직접 검색 서비스를 생성하고, 설정을 관리하며, 배포 워크플로우를 실행할 수 있는 웹 기반 어드민을 개발했습니다.
                  \item \textbf{Backend(Python/Flask)}에서는 Ansible/Digdag 실행 제어, Consul을 이용한 서비스 설정 관리 등 로직을 구현했습니다.
                  \item \textbf{Frontend(React)}에서는 운영 상태를 파악할 수 있는 SPA(Single Page Application)를 구축했습니다.
              \end{itemize}
      \end{itemize}

      %----------- [결과 및 성과] -----------%
      \textbf{결과 및 성과:}
      \begin{itemize}[leftmargin=*]
          \item \textbf{신규 검색 서비스 배포 시간 획기적 단축:} 수동으로 수일이 걸리던 신규 서비스 구축 작업을 플랫폼을 통해 수십분 -- 수시간 내에 완료할 수 있게 되었습니다.
          \item \textbf{운영 업무 리소스 절감:} 반복적인 운영 업무를 자동화하여, 검색 개발자들이 본연의 서비스 품질 개선 및 기능 개발에 집중할 수 있는 환경을 제공했습니다.
          \item \textbf{수십 개의 검색 서비스(블로그, 이미지 등)를 성공적으로 플랫폼에 전환}하여 안정적으로 운영하였습니다.
      \end{itemize}

      %----------- [사용한 기술] -----------%
      \textbf{사용한 기술:} Python, Flask, Gunicorn, Ansible, Digdag(Workflow), Consul, React
    }
  \resumeProjectListEnd

\section{\textbf{Daum/Kakao 대용량 검색 서비스 개발 및 운영}}
  \resumeProjectListStart
    \resumeProject
    {\textbf{소속}: Kakao Corp.} {2015.03 - 2019.11}

    {
      %----------- [프로젝트 개요] -----------%
      \textbf{프로젝트 개요:}
      \begin{itemize}[leftmargin=*]
          \item 일 평균 4천만 이상의 쿼리와 수백만 사용자를 감당하는 Daum 포털의 검색 서비스를 개발하고 안정적으로 운영했습니다.
          \item 평균 100ms대 응답속도, 99.99\% 이상의 서비스 가용성(SLA)이라는 매우 엄격한 목표 하에, 예측 불가능한 트래픽 속에서도 시스템 안정성을 확보하고 성능을 최적화하는 역할을 수행했습니다.
      \end{itemize}

      %----------- [핵심 역할 및 기여] -----------%
      \textbf{주요 역할 및 기여:}
      \begin{itemize}[leftmargin=*]
          \item \textbf{End-to-End 검색 컬렉션 설계 및 개발 주도}
            \begin{itemize}[leftmargin=*, label=$\circ$]
              \item 신규 검색 서비스의 요구사항 분석부터 데이터 스키마 설계, 카카오 자체 검색엔진 기반의 색인 파이프라인 구축, 검색 API 개발까지 전체 서비스 라이프사이클을 책임지고 개발했습니다.
            \end{itemize}

          \item \textbf{대규모 트래픽 환경에서의 시스템 안정성 확보 및 성능 최적화}
            \begin{itemize}[leftmargin=*, label=$\circ$]
              \item 부하 테스트를 통해 시스템의 병목 구간을 분석하고, \textbf{500 -- 1000 TPS 환경에서도 평균 100ms대의 응답 속도를 유지}하도록 최적화했습니다.
              \item 장애 시나리오 분석 및 복구 프로세스를 사전에 수립하고, 24/7 무중단 운영 환경에서 발생하는 다양한 이슈(타임아웃, 리소스 병목 등)에 대응하여 서비스 안정성을 확보했습니다.
            \end{itemize}
      \end{itemize}

      %----------- [사용한 기술] -----------%
      \textbf{담당 검색 컬렉션:} 게시판(아고라) 검색, 책 검색, 모바일앱 검색, 카카오톡 오픈채팅방 검색

      %----------- [사용한 기술] -----------%
      \textbf{사용한 기술:} Python, Golang, Java, Kakao In-house Search Engine
    }
    
  \resumeProjectListEnd
  
\end{document}
